\section{Discussion}

The measured TE routes track BF16 closely: TE-native and the
four-final-BF16-block route end 1.34\% and 0.87\% above the raw BF16 loss.
Their common probes reach 27.6K and 27.1K tokens/s/GPU. A separate
same-accelerator probe of the custom global \nvfp{} route reaches 32.5K
tokens/s/GPU; its long trajectory ends 2.69\% above BF16. TE-native \nvfp{}
and custom global \nvfp{} both use the tensor-global NVFP4 scale hierarchy;
the four-final-BF16-block route uses that hierarchy in its first 28 blocks.
The panel does not isolate the cost
of scaling from implementation and batch geometry, but it exposes the systems
constraint: a tensor-global maximum-absolute-value reduction creates a
completion dependency that the execution path must hide or pay.

Hadamard preconditioning is not uniformly beneficial across scale strategies.
The fixed-H32 \mxfp{} recipe retains 98.34\% of row-SR-only \mxfp{}
throughput while moving the raw endpoint excess from 6.53\% to 2.11\% and
validation NLL from 2.593146 to 2.422804. The fixed-H16 CTA-local recipe is
slower than CTA-local without the transform and has higher endpoint and
validation loss. These are comparisons between complete trajectories, not
experiments that isolate the effect of Hadamard preconditioning alone.

Neither tested hybrid dominates the strongest pure custom route. The
depth-wise hybrid and both operand-wise hybrids are slower and have higher
endpoint and validation loss than \mxfp{} with row-SR and fixed H32. A
controlled replay on identical saved tensors still finds lower Dgrad error for
CTA-local in all 288 tested layer--projection comparisons. The completed hybrid
trajectories do not preserve that operator-level advantage in endpoint or
validation loss.
